\uuid{48ju}
\exo7id{1663}
\titre{exo7 1663}
\auteur{ortiz}
\organisation{exo7}
\datecreate{1999-04-01}
\isIndication{false}
\isCorrection{false}
\chapitre{Réduction d'endomorphisme, polynôme annulateur}
\sousChapitre{Diagonalisation}
\cpge{mpsi-mp2i.polynomes-fractions-rationnelles.c, pcsi.polynomes.c, ptsi.polynomes.c, mpsi-mp2i.applications-lineaires.b, pcsi.applications-lineaires.b, ptsi.applications-lineaires.b, bcpst1.polynomes-reels.b, mp-mpi.reduction-endomorphismes.g, mp-mpi.reduction-endomorphismes.b, pc.reduction.a, psi.reduction-endomorphismes.a, bcpst2.valeurs-propres-vecteurs-propres.a}
\module{Algèbre}
\niveau{L2}
\difficulte{}

\contenu{
\question{
Montrer que  si le polyn\^ome minimal d'un
endomorphisme $f$ d'un $K$-espace vectoriel de
dimension finie admet une racine $\lambda\in K$
alors $\lambda$ est valeur propre de $f$.
}
}
